\chapter{性能与可维护性设计}

\section{性能设计}
LingoBridge 平台作为一款服务于跨国试点的中文自学辅助平台，其性能设计指标在物理架构选型阶段就经过了流量与计算力测算，以保障留学生在弱网环境下的听读体验。

在网络传输维度，系统设定前端静态 SPA 边缘加载延迟处于较低范围，在线课堂中的 Presence 翻页同步信令单向网络时延控制在 150 毫秒以内。

在语音合成与听读维度，由于跨国长距离请求外部 TTS 服务会产生不可控的网络延迟及调用资费，平台在后端设计了高内聚的 TTS 双层缓存调度机制。系统通过建立 L1 内存高速映射区与 L2 磁盘持久化缓冲文件系统，最大化实现拼音发音文件的就地复用。

当后端接收到 TTS 合成请求时，首先根据拼音字符串计算去噪 MD5 哈希键。若在 L1 内存缓存中命中，则直接将音频字节流返回；若在 L2 磁盘中命中，则异步将其读入内存进行 LRU 缓存置换，并返回客户端；若双层缓存均未命中，则流式调用外部合成引擎，并异步将其持久化写入磁盘与内存。该设计可实现发音接口合成开销的大幅下降。详细的缓存寻址、LRU 置换以及磁盘文件锁算法伪代码下沉到详细设计说明书（LLD）中。

在前端采集维度，针对学生进行口语作业跟读的声学硬件采集诉求，系统在前端设计了基于 Web Audio API 驱动的 MediaRecorder 拾音机制，以较低的 CPU 损耗动态收集高压缩率的 Opus 音频切片，并在内存中执行合并与物理二进制序列化，规避了直接上传大体积未压缩音频造成的网络阻塞风险。麦克风 Opus 音频帧的收集与二进制数据块序列化生成逻辑详细伪代码下沉到详细设计（LLD）中。

\section{可扩展性设计}
本平台的可扩展性设计核心在于将数据持久层进行抽象，实现了松耦合的双模数据底座拔插架构。

后端开发团队统一使用标准的 Repository 仓库模式声明外部业务数据存取契约。当后续试点规模扩大，需将试点小班级从极简冷启动的 JSON 存储迁移到高性能的关系型 PostgreSQL 数据库时，控制路由层及业务逻辑层代码无需做任何物理重构，仅需将系统的运行时环境变量 \texttt{DB\_MODE} 修改为 \texttt{postgres}，仓库模块就会自动将数据交互管道切换为向物理连接池分发高效的结构化 SQL，在工程上达成了无缝平滑扩容的目标。

\section{可维护性设计}
为了保障平台在后续试点和生产运维阶段的低维护成本，系统在模块分层、运维监控及代码构建维度制定了严格的规范。

在模块物理划分上，前端 UI 展示组件、接口调用层、路由控制层、数据访问仓库层有着严格的上下游依赖单向路径，坚决杜绝任何反向跨层引用，极大地降低了后期业务变更造成代码污染的风险。

在生产运维维度，后端 Node.js 信令服务通过成熟的进程管理器 PM2 进行常驻守护，配置了完备的多核 CPU 负载均衡与崩溃物理自启机制。一旦检测到后端进程因异常阻塞，PM2 会在微毫秒内安全执行热自启，并通过内置日志输出管道对错误信息进行本地截获，确保整个平台在零运维专家值守的情境下依然能够实现常驻运行。

\section{兼容性设计}
平台前端在兼容性层面上，摒弃了非标准的前端组件与冷门多媒体格式，全面对齐国际 W3C 标准规范。

在声学采集和媒体播放维度，前端统一使用高兼容性的 Web Audio API 进行物理采样，并以各大主流浏览器均支持的 WebM 与 Opus 编码作为多媒体传输容器，确保留学生使用 Chrome、Edge、Firefox 甚至 Safari 浏览器，均能获得一致的课件白板同步、音视频互动以及发音纠音评测体验。
